\SourceLeaf{291}%
\[
\begin{aligned}
\sin(Q+q)&=\sin Q\cos\tfrac{1}{2}(h-\delta)+\cos Q\sin\tfrac{1}{2}(h-\delta)\cos\chi\\
\sin(Q-q)&=\sin Q\cos\tfrac{1}{2}(h+\delta)-\cos Q\sin\tfrac{1}{2}(h+\delta)\cos\chi
\end{aligned}
\]
und ihre Summe mit \(2\cos Q\) dividirt
\[
\operatorname{tang}Q\mathord{.}(\cos q-\cos\tfrac{1}{2}h\mathord{.}\cos\tfrac{1}{2}\delta)=-\cos\tfrac{1}{2}h\sin\tfrac{1}{2}\delta\cos\chi
\]
\indent Da nun offenbar \(\cos q-\cos\tfrac{1}{2}h\mathord{.}\cos\tfrac{1}{2}\delta\) eine Grösse zweiter Ordnung ist, so wird auch \(\sin\tfrac{1}{2}\delta\cos\chi\), und \(\delta\cos\chi\) von dieser Ordnung sein, mithin, da \(\lambda\) den Factor \(\cos\chi^2\) implicirt, \(\lambda hh\delta\) von der vierten, und \(\lambda h\delta\delta\) von der fünften Ordnung; hiedurch ist also die Weglassung dieser Glieder gerechtfertigt.

Das Endresultat dieser Entwicklung ist demnach, unter der angegebenen Voraussetzung, in folgenden Formeln enthalten, wo anstatt der Tangenten von \(\psi^0,\psi'\) die Bögen selbst geschrieben sind:
\[
\begin{aligned}
\psi^0&=-\frac{ee\cos P\sin P\sin\chi\cos\chi^2h^3}{12\cos\varphi\cos\theta}\\
\psi'&=+\frac{ee\cos P\sin P\sin\chi\cos\chi^2h^3}{12\cos\varphi\cos\theta}
\end{aligned}
\]
\Article{16}
Die Berechnung des Dreieckssystems auf der Kugel zerfällt in drei Hauptstücke:
\begin{list}{}{\setlength{\leftmargin}{1.8em}\setlength{\labelwidth}{1.4em}\setlength{\labelsep}{.4em}\setlength{\itemsep}{0pt}}
\item[1)] die Ausgleichung der Winkel nach allen den Bedingungsgleichungen, welche die Beschaffenheit des Systems darbietet.
\item[2)] die Berechnung der sämmtlichen Dreiecksseiten.
\item[3)] \begingroup\rightskip=0pt plus 1em
die Bestimmung der Längen und Breiten der Dreieckspunkte, in Verbindung mit der Orientirung der von jedem derselben ausgehenden Dreiecksseiten.\par\endgroup
\end{list}
\hspace*{\parindent}Die Verwandlung der Längen und Breiten auf der Kugel in die wahren Längen und Breiten auf dem Sphaeroid geschieht dann für die Längen durch die Division mit dem constanten Divisor \(\alpha\), für die Breiten vermittelst der hier beigefügten Hülfstafel, oder einer andern auf ähnliche Weise besonders construirten, wenn man einen andern Normal-Parallelkreis zu wählen Ursache hat.

Mit Übergehung der beiden ersten auf bekannten Gründen beruhenden Geschäfte füge ich hier noch einiges in Beziehung auf das dritte bei, welches sich auf die Auflösung der Aufgabe reducirt\footnote{Da diese Aufgabe hier wie eine für sich bestehende betrachtet wird, so können ohne Nachtheil einige Buchstaben hier in anderer Bedeutung als oben gebraucht werden.}: aus der in Bogentheilen ausgedrückten
